\documentclass[11pt]{article}

\usepackage{amsmath}
\usepackage{amsfonts}
\usepackage{graphicx}
\usepackage{hyperref}
\usepackage{physics}
\usepackage{bm}

\title{Mathematical Foundations of the MedVision AI Pipeline}
\author{Yann Smatti}
\date{2026}

\begin{document}

\maketitle

\tableofcontents

\newpage

\section{Introduction}

MedVision AI is a computer vision pipeline designed to analyze medical images
using deep learning techniques.

The system includes the following components:

\begin{itemize}
\item Image preprocessing
\item Convolutional feature extraction
\item Transfer learning
\item Model optimization
\item Experiment tracking with MLflow
\end{itemize}

This document provides a mathematical description of the algorithms
used in the project.

\section{Mathematical Representation of Images}

A digital image can be modeled as a discrete tensor.

Let

\[
I : \Omega \subset \mathbb{Z}^2 \rightarrow \mathbb{R}^C
\]

where

\begin{itemize}
\item $C$ is the number of channels
\item $\Omega = \{1,\dots,H\} \times \{1,\dots,W\}$
\end{itemize}

Thus

\[
I \in \mathbb{R}^{H \times W \times C}
\]

For RGB images $C = 3$.

\subsection{Image as a vector in high dimensional space}

An image can also be represented as a vector

\[
x \in \mathbb{R}^{HWC}
\]

For example:

\[
224 \times 224 \times 3 = 150\,528
\]

Thus a single image lives in a
150k-dimensional space.

Computer vision can therefore be interpreted
as learning functions

\[
f : \mathbb{R}^{150000} \rightarrow \mathbb{R}
\]

which explains why deep neural networks are necessary.

\section{Discrete Convolution}

The convolution operation is central to CNNs.

Given an image $I$ and kernel $K$:

\[
(I*K)(i,j) =
\sum_{m=-k}^{k}
\sum_{n=-k}^{k}
I(i-m,j-n)K(m,n)
\]

\subsection{Connection with Fourier analysis}

Convolution in the spatial domain corresponds
to multiplication in the Fourier domain.

\[
\mathcal{F}(f * g) =
\mathcal{F}(f)\mathcal{F}(g)
\]

This explains why convolution acts as a filter.

\section{Feature Extraction}

A convolutional layer computes

\[
Y_{i,j,f} =
\sum_{c=1}^{C}
\sum_{m=0}^{k-1}
\sum_{n=0}^{k-1}
W_{m,n,c,f}X_{i+m,j+n,c}
\]

Important properties:

\begin{itemize}
\item local receptive field
\item weight sharing
\item translation invariance
\end{itemize}

\section{Gradient Based Optimization}

The model parameters $\theta$ are optimized by minimizing
a loss function $L(\theta)$.

The update rule is

\[
\theta_{t+1} =
\theta_t - \eta \nabla L(\theta_t)
\]

\subsection{Backpropagation}

Using the chain rule

\[
\frac{\partial L}{\partial x}
=
\frac{\partial L}{\partial y}
\frac{\partial y}{\partial x}
\]

Gradients propagate through the network
layer by layer.

\section{Transfer Learning}

The network is decomposed into

\[
f(x) = g(h(x))
\]

where

\begin{itemize}
\item $h(x)$ extracts generic visual features
\item $g(x)$ performs task specific classification
\end{itemize}

This dramatically reduces the required dataset size.

\section{Explainability}

Grad-CAM generates visual explanations.

\[
\alpha_k^c =
\frac{1}{Z}
\sum_{i,j}
\frac{\partial y^c}{\partial A_{i,j}^k}
\]

\[
L^c =
ReLU
\left(
\sum_k \alpha_k^c A^k
\right)
\]

This highlights regions influencing the decision.

\end{document}